\section{Discussion and Limitations}
\label{sec:discussion}

This section discusses the principal limitations of the work and the
directions we see as most promising for follow-up.

\paragraph{Single-container warm pools.} Lemma~\ref{lemma:tradeoff} is
stated for a single warm container. Real FaaS platforms maintain
warm \emph{pools} of size $n \ge 1$. The multi-container generalization
changes the survival probability $\Pr[X > T_w]$ to a function of the
steady-state distribution of an M/M/$n$-style queue, and scales idle
energy with $\mathbb{E}[n - \mathrm{busy}]$. We conjecture an
analogous threshold structure carries over, but the exact dependence
on $n$ involves the (non-closed-form) steady-state occupancy
distribution and remains future work. The scheduler currently applies
the single-container lemma to each incremental container in the pool,
which is the right limit-case behaviour but does not exploit
joint-pool dynamics. The empirical impact in the regimes we tested
is small.

\paragraph{Reach of the dichotomy on FaaS parameters.} As we
acknowledge in \S\ref{sec:tradeoff:dichotomy}, the lemma's first
regime ($\beta \le \beta_{\mathrm{crit}}$, warm-in-L wins for all
$\lambda > 0$) is effectively unreachable for realistic FaaS function
profiles, where $\beta = P_i T_w / (P_a \tau_e) \in [50, 200]$ and
$\beta_{\mathrm{crit}} = (r-1)(1+\alpha) \le 47$ even for extreme
carbon ratios. The practical content of the lemma for FaaS
scheduling is therefore the second regime: the closed-form
computation of $\lambda^*$ that the scheduler uses to gate region
admission. The dichotomy formulation has theoretical value (it is
the natural starting point for the proof) and is reachable for
adjacent computing classes (long-running idle keep-alives, edge
containers with short TTLs), but the FaaS-specific practical
contribution is the $\lambda^*$ characterization, not the
classification.

\paragraph{No demonstrated dominance over pure spatial routing.}
A careful read of the empirical results in
\S\ref{sec:evaluation:realtopology} reveals that \GreenFaaS{} does
not strictly outperform the simpler Spatial baseline on any
real-carbon topology we tested. The two schedulers stay within 1.3
percentage points across five real-data topologies, and Spatial in
fact edges \GreenFaaS{} on real ENTSO-E coal-belt data
(\S\ref{sec:evaluation:realtopology}). \GreenFaaS's distinct value
relative to Spatial is therefore not ``better carbon reduction'' but
``topology-robust scheduler that does not require operator tuning
for low-carbon-refuge presence, includes a principled cold-start
trade-off lemma, and preserves do-no-harm across a fine-grained
range of low-opportunity workloads
(\S\ref{sec:evaluation:fine-donoharm})''. We believe this is a more
defensible and practically useful claim than ``always wins''. The
underlying point is that on real-world grid carbon data with
realistic FaaS arrival patterns, the spatial axis already captures
most of the carbon-aware savings available; the temporal-shift
contribution is necessary for robustness but does not unlock large
additional headroom on top of well-tuned spatial routing.

\paragraph{Provable online algorithms on FaaS.} Our
\S\ref{sec:evaluation:lechowicz} comparison against a FaaS-adapted
port of \citet{lechowicz2024online}'s provably optimal one-way trading
algorithm produced a striking empirical finding: the provably-optimal
algorithm performs \emph{worst} of all baselines on FaaS workloads,
including worse than the heuristic Wait-Awhile. The mechanism is the
non-separable cost structure --- warm-pool idle energy couples
decisions across invocations through state that Lechowicz's
competitive analysis does not see, so the analytically-optimal
threshold is empirically wrong on FaaS. We have not attempted to
extend the competitive-ratio proof to the warm-pool-coupled cost
structure; the existing framework's separable-per-job-cost assumption
is fundamental to the analysis and a FaaS-native online algorithm
with competitive guarantees over the coupled problem remains the
most interesting open theoretical extension of this work.

\paragraph{Workload-trace authenticity.} The sensitivity sweeps of
\S\ref{sec:evaluation:variability}--\S\ref{sec:evaluation:intensity}
use a schema-faithful synthetic workload calibrated to the Azure
Functions 2019 and 2021 characterizations
\citep{shahrad2020serverless,zhang2021faster,azuretrace}, because
parameter variation across multiple axes is the purpose of those
experiments. The headline and topology experiments
(\S\ref{sec:evaluation:headline},
\S\ref{sec:evaluation:realboth},
\S\ref{sec:evaluation:realtopology})
report both the synthetic and the
fully-real configurations: real Azure 2021 per-invocation trace
($1.98 \times 10^6$ invocations over 14 days, $\sim$81k in a 24h
slice) paired with real ElectricityMaps carbon-intensity data from
the same January 2021 window. Time-aligned fully-real results are
in \S\ref{sec:evaluation:realboth}, Table~\ref{tab:realboth_headline}.
The qualitative findings of the synthetic experiments survive
unmodified on fully-real data; the quantitative magnitudes of the
batch-scheduler failure modes are softer on real data (3--5\% loss vs
FIFO instead of 200--300\%), reflecting the real workload's higher
baseline cold-start rate. We note this softening transparently in
\S\ref{sec:evaluation:realboth}.

\paragraph{Non-Poisson arrivals.} The lemma assumes Poisson
inter-arrival times, empirically violated by real production
workloads \citep{shahrad2020serverless}. The evaluation on real LWA
carbon traces with a schema-faithful bursty workload
(\S\ref{sec:evaluation:realcarbon}) shows that the scheduler still
achieves 64\% reduction vs FIFO, within 0.5 percentage points of
pure spatial routing; we interpret this as evidence that the lemma's
\emph{qualitative} prescription is robust to non-Poisson arrivals
even though the \emph{exact} break-even rate $\lambda^*$ shifts. A
formal extension to general renewal processes is open.

\paragraph{PUE heterogeneity.} The lemma assumes similar PUE across
regions. Real data centres span a non-trivial range (PUE 1.1 in
hyperscale Nordic sites; PUE 1.5 in older facilities). The lemma
extends by replacing $r = C_H/C_L$ with
$r' = (C_H \cdot \mathrm{PUE}_H)/(C_L \cdot \mathrm{PUE}_L)$, but our
evaluation uses a uniform $\mathrm{PUE} = 1.2$, close to the Uptime
Institute's reported global-average data-centre PUE in recent annual
surveys (which has plateaued near $1.5$ industry-wide but is closer
to $1.1$--$1.2$ for modern hyperscale facilities of the kind that run
multi-region FaaS). A PUE-heterogeneous
evaluation with measured per-site values is a natural follow-up.

\paragraph{Robustness to power-model assumptions.} Our simulator uses
$P_a = 4$\,W active and $P_i = 0.3$\,W idle as defaults, calibrated
to typical container-level draws. Because production-measured power
(via Kepler, \S\ref{sec:deployment:power}) may differ, we verified
that the scheduler ranking is invariant to the power model. Scaling
$P_a$ and $P_i$ by factors in $\{0.5, 1.0, 1.5\}$, including the
asymmetric perturbations $(P_a {\times} 0.5, P_i {\times} 1.5)$ and
$(P_a {\times} 1.5, P_i {\times} 0.5)$, leaves the ordering
Spatial $<$ \GreenFaaS{} $<$ GreenFaaS-v1 $<$ FIFO $<$ Wait-Awhile
unchanged in every case (real Azure 2021 + real EM Jan 2021, 5-region
topology). The absolute carbon scales with the power assumptions, as
expected, but the relative scheduler performance --- the paper's
actual claim --- does not depend on the specific $P_a/P_i$ values.

\paragraph{Cold-start-time sensitivity.} The simulator treats the
cold-start duration $\tau_c$ as a constant per-function property
($\sim 1$\,s in our catalog). In production, $\tau_c$ varies with
language runtime and image-pull state, ranging from sub-second to
several seconds, and concurrent cold starts can contend for I/O. Since
$\alpha = \tau_c/\tau_e$ enters the Lemma~\ref{lemma:tradeoff}
threshold, larger $\tau_c$ raises $\beta_{\mathrm{crit}}$ and shifts
$\lambda^*$, making the gate \emph{more} willing to keep a warm
container. We did not sweep $\tau_c$ explicitly; a $0.5$--$5$\,s span
(a $10\times$ range) changes $\lambda^*$ monotonically but does not
reverse the qualitative gate decisions for our workload's arrival
rates, because the spatial carbon ratio dominates the gate in the
multi-region topologies where savings occur.

\paragraph{Mean-carbon estimation transient.} The idle-energy
correction that enables do-no-harm uses $\bar{C}_r$, the long-run mean
carbon intensity per region. During the first hours of a fresh
deployment a running estimate of $\bar{C}_r$ may be biased, weakening
the correction and potentially admitting an unprofitable deferral. A
production deployment should seed $\bar{C}_r$ with a historical annual
mean (e.g. from ElectricityMaps) rather than a cold running average,
and log per-region $\bar{C}_r$ estimates alongside deferral counts for
post-hoc auditing.

\paragraph{No competitive-ratio bound.}
As discussed in \S\ref{sec:problem:online}, we make no claim of
competitive optimality. The non-separable cost structure (idle energy
couples decisions across invocations through the warm-pool state)
together with the adversarial-perturbation argument rule out a finite
competitive ratio in the worst case. The Lechowicz et al.\
double-threshold framework \citep{lechowicz2024online} gives optimal
ratios for a strictly simpler model and is a natural starting point
for a future theoretical extension. The principal obstacle is the
warm-pool idle-energy term, which has no analogue in pause/resume. We
view this as the most interesting open theoretical question raised
by the paper.

\paragraph{Embodied carbon.}
\label{sec:discussion:embodied}
\GreenFaaS{} optimizes \emph{operational} carbon and ignores
\emph{embodied} carbon. EcoLife \citep{jiang2024ecolife} addresses
embodied carbon by routing functions to older, already-amortised
hardware. This is an orthogonal axis to the spatial/temporal grid-
carbon axes we exploit; a natural composition would route an
invocation first by \GreenFaaS's spatial/temporal logic to a
region/time, and then within that region by EcoLife's
hardware-generation logic. We have not built this composition; it is
the most concrete next-paper-sized direction this work suggests.

\paragraph{Real-data validation.} Both axes of our evaluation are
validated on real data. The grid-carbon side uses real ENTSO-E /
CAISO 2020 data for four regions
(\S\ref{sec:evaluation:realcarbon}, \S\ref{sec:evaluation:realtopology})
and real ElectricityMaps January 2021 data for five regions
(\S\ref{sec:evaluation:realboth}). The workload side is validated on
the real Azure Functions 2021 per-invocation trace
($1.98 \times 10^6$ invocations over 14 days), paired with
time-aligned ElectricityMaps carbon data from the same January 2021
window, in the fully-real experiments of
\S\ref{sec:evaluation:realboth}. The controlled sensitivity sweeps
(\S\ref{sec:evaluation:sla}--\S\ref{sec:evaluation:intensity}) use a
schema-faithful synthetic workload because parameter variation across
multiple axes is their purpose, but the headline ordering is
confirmed on the fully-real data: the Wait-Awhile and Lechowicz
failure modes, the \GreenFaaS/Spatial near-match, the v1 ablation
gap, and the do-no-harm property all reproduce, with the
batch-scheduler failure magnitudes softer on real data (3--5\% vs
FIFO) than on synthetic (up to $-306\%$), reflecting the real
workload's higher baseline cold-start rate.

\paragraph{Real-trace scope and LLM workloads.} The Azure traces date
from 2019 and 2021. The serverless mix has evolved since then, most
notably with LLM inference functions, which have substantially longer
durations and higher per-invocation energy. Our latency-class
heuristics would classify LLM inference as background, which is
technically correct under our definition but understates the
schedulable opportunity for interactive LLM applications such as
chatbots. We expect the qualitative findings to carry over ---
\GreenFaaS's do-no-harm property is workload-independent --- but
absolute numbers would shift on a contemporary LLM-heavy trace.

\paragraph{Production-system validation.} \GreenFaaS{} is evaluated
entirely in simulation. The Lemma~\ref{lemma:tradeoff} break-even
computation is cheap enough ($<100\,\mu$s per call) that integration
into a production scheduler such as Knative or AWS Lambda is
plausible; the closest production point of comparison, GreenCourier
\citep{chadha2023greencourier}, is a Kubernetes plugin with
$\sim$24\,ms binding-latency overhead above the default scheduler. A
multi-region Knative deployment with measured grid carbon would be
the strongest possible confirmation of the simulator results.

\paragraph{Summary.} We interpret these limitations not as caveats
that weaken the contributions but as \emph{scope statements} that
delimit what we have shown. The closed-form trade-off
characterization, the SLA-tiered algorithm, the do-no-harm empirical
property, and the forecast-invariance result are robust within the
stated assumptions. The interesting next steps are the
multi-container generalization, the embodied-carbon composition with
EcoLife, and the production-system validation --- each a tractable
follow-up rather than a fundamental obstacle.
